\chapter{Introducción}
\label{ch:introduccion}

\section{Propósito del documento}

El presente documento corresponde a la \textbf{Revisión Crítica de Diseño (CDR)}
del \textbf{Subsystema de Super-Resolución basada en Inteligencia Artificial
(SAI)} del proyecto CubeSat INTISAT, desarrollado en el Instituto de
Radioastronomía (INRAS) de la Pontificia Universidad Católica del Perú (PUCP).

El objetivo de este CDR es documentar el diseño, la arquitectura, los
requerimientos, la estrategia de verificación y los riesgos asociados al
subsistema SAI, previo a su integración con el \textit{CMOS Microscopy Payload}.
Este documento sigue lineamientos de ingeniería de sistemas ECSS/NASA, adaptados
al contexto académico del proyecto INTISAT.

\section{Rol del subsistema SAI dentro de INTISAT}

El subsistema SAI constituye un \textbf{módulo de postprocesamiento
computacional}, cuyo rol es mejorar la \textbf{interpretabilidad visual} de las
imágenes adquiridas por el CMOS Microscopy Payload (CMOSM).

Es importante destacar que el SAI:
\begin{itemize}
    \item No es un instrumento de adquisición.
    \item No genera información física nueva.
    \item No modifica los datos originales capturados por el sensor.
\end{itemize}

Su función principal es actuar como una etapa posterior al proceso de captura,
aplicando técnicas de aprendizaje profundo para realzar contraste, nitidez y
continuidad espacial de estructuras ya presentes en la imagen base.

\section{Motivación para el uso de super--resolución}

Las imágenes obtenidas mediante microscopía sin lentes en configuración de
contacto presentan limitaciones inherentes relacionadas con:
\begin{itemize}
    \item tamaño de píxel del sensor,
    \item ruido electrónico,
    \item difracción y dispersión cercana,
    \item restricciones de iluminación.
\end{itemize}

En este contexto, la super--resolución basada en aprendizaje profundo se evalúa
como una herramienta complementaria para:
\begin{itemize}
    \item mejorar la legibilidad visual de estructuras celulares,
    \item facilitar la interpretación humana de las imágenes,
    \item apoyar procesos de análisis cualitativo.
\end{itemize}

El uso de técnicas de super--resolución se plantea de manera conservadora,
reconociendo explícitamente sus limitaciones y evitando interpretaciones
científicas basadas exclusivamente en imágenes procesadas.

\section{Alcance del subsistema SAI}

El alcance del subsistema SAI en la presente fase de desarrollo se limita a:

\begin{itemize}
    \item Procesamiento de imágenes previamente adquiridas por el CMOSM.
    \item Ejecución de modelos de super--resolución en entorno de laboratorio o
    en tierra.
    \item Evaluación experimental del impacto visual del procesamiento.
\end{itemize}

Quedan explícitamente fuera del alcance de este CDR:
\begin{itemize}
    \item La inferencia en tiempo real a bordo del satélite.
    \item El uso de IA como fuente primaria de datos científicos.
    \item La certificación del subsistema para operación en entorno espacial.
\end{itemize}

\section{Relación con el CMOS Microscopy Payload}

El subsistema SAI se integra funcionalmente como una extensión del flujo de datos
del CMOSM. La frontera entre ambos subsistemas se define de la siguiente manera:

\begin{itemize}
    \item El CMOSM es responsable de la formación física de la imagen y del
    registro de metadatos.
    \item El SAI recibe imágenes ya calibradas y las procesa sin alterar los
    datos originales.
\end{itemize}

Esta separación clara de responsabilidades permite mantener la trazabilidad de
los datos y evita dependencias críticas entre adquisición y procesamiento.

\section{Estructura del documento}

El presente documento se organiza de la siguiente manera:

\begin{itemize}
    \item \textbf{Capítulo~\ref{ch:introduccion}}: Introducción, rol y alcance del
    subsistema SAI.
    \item \textbf{Capítulo~2}: Requerimientos funcionales, computacionales y de
    interfaz del subsistema.
    \item \textbf{Capítulo~3}: Arquitectura y diseño del pipeline de
    super--resolución.
    \item \textbf{Capítulo~4}: Verificación y evaluación experimental.
    \item \textbf{Capítulo~5}: Análisis de riesgos y limitaciones.
    \item \textbf{Capítulo~6}: Conclusiones y trabajo futuro.
\end{itemize}
